\scanpage{20}
\deleted{3) Si $\underline E$ est \deleted{\unsure{de}} de t.f. \inserted{et est quotient d'un module loc libre (\unsure{sans op.!})}, il existe un
$\unsure{\mathcal O|\mathcal E|\underline{O}}_X$-$G$-Module \deleted{\unsure{qu.}} \uline{loc. libre} \uline{de t.f.} $\underline F$,
et un épimorphisme de $\mathcal O_X$-$G$-Modules
$\underline F\longrightarrow\underline E$.}

3) \deleted{Si $S$ est \unsure{p.} admet un faisceau inv.}
\inserted{qu. cpt \deleted{\unsure{qu. sép.}} qu. sép. et si tt module quasi coh.}
de t.f. sur $S$ est $\simeq$ quotient d'un loc. libre de t.f.,
\deleted{ample,} alors $\underline E$ est $\unsure{\simeq|\cong}$ quotient d'un
$\underline O_X$-$G$-Module loc. libre de type fini
de la forme $f^*(H)(-n)$, où

\uline{Dém} Il suffit de prouver 1) et 3). Posons
(\unsure{où|ici} $f:X\unsure{-|\to}S$)
\[
\underline{\Gamma}_{\deleted{\unsure{*}}}^{\unsure{\bullet|\cdot}}(E)=\coprod_n f_*(E(n))
\]
Alors grâce au fait que $G$ est plat sur $S$
et $f$ qu. cpt et (qu) séparé, on trouve que
$\Gamma^\bullet(E)$ est un Module gradué sur $S$ sur
lequel $G$ opère. Soit
\[
\underline{\Gamma}_{\deleted{*}}^\bullet(\mathcal O_X)=\deleted{\unsure{A_*|A^\bullet}}\ A^\bullet
\]
c'est une Algèbre graduée sur \unsure{lesquels|laquelle}
$G$ opère. $E\unsure{\leadsto|\sim|\mapsto}\underline{\Gamma}^{\cdot}(E)$ est un foncteur en
$E$, à valeurs dans les \deleted{\unsure{objets}}
$A^{\cdot}$-$G$-Modules gradués. Ce foncteur
est exact à gauche. On a un foncteur
(\uline{exact})
% archive sheet number: 19
